Automated vulnerability detection in source code is critical for software
security. Existing approaches fall into two camps: image-based methods
(e.g., VulCNN, VulGAI) that convert code into visual representations for
CNN classification — fast but losing structural relationships — and graph-based
methods (e.g., ReGVD, iGnnVD) that apply graph neural networks (GNNs) over
code graphs — accurate but ignoring visual and semantic signals. No prior work
unifies all three modalities for general source-code vulnerability detection.

We present \textbf{VulGCL}, a multimodal vulnerability detection framework that
jointly learns from three complementary views of a code function: (1)~a
\emph{graph branch} that applies a GNN over the Program Dependency Graph (PDG)
to capture structural relationships; (2)~an \emph{image branch} that converts
the PDG into a three-channel image via node centrality encoding and applies a
CNN to capture visual patterns; and (3)~an \emph{LLM branch} that fine-tunes
CodeBERT to capture the semantic meaning of code. The three representations are
fused and passed to a vulnerability classifier.

We additionally propose \textbf{VulGCL-Lite}, a quantized version of VulGCL
that runs inference on a Raspberry Pi 4B, enabling real-time vulnerability
scanning on resource-constrained IoT and embedded devices.

Experiments on the Devign and BigVul benchmarks show that VulGCL outperforms
all single-modality baselines and achieves state-of-the-art F1 of XX.X\% on
Devign and XX.X\% on BigVul. VulGCL-Lite achieves XX.X\% of full-model
accuracy at XX\,ms per sample on Raspberry Pi 4B.
% ↑ Fill in actual numbers after experiments
